% Secondary A.2(b) proof by the Part C agent.  Root may integrate this fragment.
% The external geometric inputs are infinite-grafting uniformization,
% the collar lemma, parameter dependence, and the weighted reduced-modulus theorem.
\subsection{The exact radius image and uniqueness of canonical caps}
\label{sec:A2b-secondary}

Let $Y$ be a connected compact marked Riemann surface of genus $g$, with
distinct marked points $x_1,\ldots,x_n$, where $2g-2+n>0$.
Choose a nonempty set $I$ of points that are to receive genuine caps;
the other marked points remain cusps. Fix a nonzero tangent coordinate
at each $x_i$, $i\in I$, so that conformal radii are positive numbers.
For $s=(s_i)_{i\in I}\in(0,\infty)^I$, put
\[
 \lambda_i=\sqrt{s_i},\qquad L_i=2\pi\lambda_i.
\]
Infinite-grafting uniformization gives the unique marked hyperbolic surface
with geodesic boundaries of these lengths whose infinite grafting has
conformal compactification $Y$. Denote its cap discs by $D_i(s)$, their
conformal radii by $r_i(s)$, and set $F_i(s)=\log r_i(s)$.
This use of the grafting theorem fixes the compact conformal surface,
as well as the indicated zero boundary lengths.

There is a related variational formula in the literature. In the sphere
setting treated by F\i rat, the on-shell action satisfies
$\partial S_{\mathrm{HJ}}/\partial\lambda_i=2\lambda_i\log r_i$
\cite[Appendix B, equation (B.1)]{AFirat2023}; after the change
$s_i=\lambda_i^2$, its derivative is $\log r_i$. We do not infer strict
convexity or the global range from that identity. The moving-seam energy
calculation below supplies these properties for every stable genus in
the present scope. We claim no priority for the variational viewpoint.

\begin{lemma}[Proved: regularity needed for variation]\label{A2b-regularity}
Locally in $(0,\infty)^I$, the canonical seams, their conformal cap maps,
and the hyperbolic and flat parts of the grafted metric depend
real-analytically on $s$. The logarithm of the metric density is $C^{1,1}$
across each seam and is smooth on either side. Its first parameter
variation is continuous across the seam and has the one-sided normal
derivatives used below.
\end{lemma}
\begin{proof}
The real-analytic inverse in the finite strata of the infinite-grafting
theorem gives a real-analytic family of marked hyperbolic surfaces, with
fixed cusp ends. Choose analytic Fenchel--Nielsen representatives locally.
Geodesic boundary collars and their normalized arc-length parameters
depend analytically on these representatives. Appending the standard
cylinders therefore gives an analytic family of complex charts, although
the metric, expressed in a chart crossing a seam, is only $C^{1,1}$.
The marked conformal isomorphisms of the compactifications with the fixed
$Y$ can be chosen real-analytically: apply parameter dependence for the
normalized Beltrami equation in these charts
\cite[\S5.7, Theorems 10--11, pp.~402--403]{AhlforsBers1960A2b}.
One can make the descent explicit on the universal cover of the marked
punctured surface. The normalized Beltrami solutions and the conjugated
deck transformations depend analytically on the parameters. Because
the marked conformal structure is fixed, these deck groups are conjugate
to the fixed group. Normalize that conjugacy by three distinct boundary
fixed points of deck transformations; its coefficients are analytic.
The resulting maps descend to the desired marked isomorphisms. In the
stable marked situation there is no nonidentity holomorphic automorphism
isotopic to the identity relative to the marked points. The resulting
maps and their parameter derivatives are smooth in the spatial variable
on a chart crossing the seam: its Beltrami coefficient is smooth, so
interior elliptic estimates applied to the Beltrami equation and its
parameter derivatives upgrade parameter dependence from the spaces in
that theorem to every local differentiability order. This proves the
claimed parameter regularity.

More explicitly, at a fixed seam use $\zeta=x+iy$, where $x$ has period
$2\pi\lambda_i$, $y>0$ is the flat side, and the initial metric is
\[
 g=\begin{cases}\sec^2y\,|d\zeta|^2,&y<0,\\
                  |d\zeta|^2,&y>0.
   \end{cases}
\]
The two logarithmic densities and their first derivatives agree at
$y=0$. Their second normal derivatives have finite one-sided values.
Differentiating their equality along the analytic moving seam gives
precisely the continuity and jump conditions calculated below; no
second derivative across the seam is assumed.
\end{proof}

\begin{theorem}[Proved: strict monotonicity of logarithmic cap radii]
\label{A2b-energy}
The Jacobian $dF/ds$ is symmetric positive definite everywhere.
Consequently $F$ is the gradient of a strictly convex real-analytic
function $S$ on $(0,\infty)^I$, and is injective. In particular an
admissible canonical cap system is determined by its conformal radii.
\end{theorem}
\begin{proof}
Take a tangent variation, and write
\[
 g_t=e^{2u_t}|dz|^2,\qquad h=\dot u_0,\qquad
 c_i=\dot\lambda_i/\lambda_i.
\]
The ratio defining $h$ is independent of the fixed local coordinate.
Let $X$ denote the hyperbolic core, including its cusp ends, and $C_i$
the flat cap cylinder. Differentiating the curvature equations gives
\[
 (\Delta_g-2)h=0\quad\hbox{on }X,\qquad
 \Delta_g h=0\quad\hbox{on }C_i.
\]
Use the collar coordinate of Lemma~\ref{A2b-regularity}, so the initial
cap coordinate is $w_i=\exp(i\zeta/\lambda_i)$.
Normalize the varying cap coordinate by its positive derivative in the
chosen tangent coordinate. The function
\[
 k_i=\operatorname{Re}\left.\frac{d}{dt}\right|_0\log w_{i,t},
 \qquad V_i=\lambda_i k_i
\]
is single-valued and harmonic on $C_i$. Since
$w_{i,t}(z)=z/r_i(t)+O(z^2)$ at its puncture,
\[
 V_i=-\lambda_i\dot F_i+O(e^{-y/\lambda_i}),\qquad
 h=c_i+O(e^{-y/\lambda_i})\quad(y\longrightarrow\infty).
\]
Differentiating
$g_t=\lambda_i(t)^2|dw_{i,t}/w_{i,t}|^2$ gives
\begin{equation}\label{A2b-cap-variation}
 h=c_i-\partial_y V_i.
\end{equation}
The equation $|w_{i,t}|=1$ shows that $V_i(x,0)$ is the normal velocity
of the seam in the direction from the core towards the flat cylinder.
The initial logarithmic density is $-\log\cos y$ on the core and $0$
on the flat side. Differentiating equality of its first derivatives
along the moving seam gives
\begin{equation}\label{A2b-jump}
 (\partial_y h)_+-(\partial_y h)_-=V_i.
\end{equation}
The sign follows also directly from the second derivatives $0$ and $1$
of those two logarithmic densities. Here $+$ denotes the flat side.

All the following integrals converge. The flat-end estimates above give
exponential decay of derivatives. At a fixed cusp, writing its canonical
coordinate as $q_t(z)=a_tz+O(z^2)$ in
$|dq_t/q_t|^2/(\log|q_t|)^2$ shows that
$h=O(1/|\log|z||)$ and that the truncated boundary terms tend to zero.
Green's identity, first on truncated pieces and then in the limit, yields
\begin{align*}
 E_h&:=\int_X\bigl(|\nabla h|^2+2h^2\bigr)dA
       +\sum_i\int_{C_i}|\nabla h|^2dA
       =-\sum_i\int_{\Gamma_i}hV_i\,dx,\\
 E_V&:=\sum_i\int_{C_i}|\nabla V_i|^2dA
       =\sum_i\int_{\Gamma_i}hV_i\,dx
          -\sum_i c_i\int_{\Gamma_i}V_i\,dx.
\end{align*}
In the first identity the core outward normal is $+\partial_y$ and
the cylinder outward normal is $-\partial_y$; use
\eqref{A2b-jump}. In the second use
\eqref{A2b-cap-variation}. The circular mean of a bounded harmonic
function on a half-cylinder is constant, so
\[
 \int_{\Gamma_i}V_i\,dx=-2\pi\lambda_i^2\dot F_i.
\]
Adding the identities proves the exact formula
\begin{equation}\label{A2b-positive-energy}
 \boxed{\quad
  \pi\sum_i\dot s_i\dot F_i
       =E_h+E_V.\quad}
\end{equation}
The right side is positive for a nonzero tangent vector: if it is zero,
$h=0$ on the nonempty hyperbolic core; continuity and the bounded
harmonic Dirichlet problem on each cylinder then give $h=0$ there,
and hence every $c_i=0$.

For two tangent variations $a,b$, apply the same integrations by parts
bilinearly. Their left side is
$\pi\sum_i\dot s_i^{,a}\dot F_i^{,b}$ and their right side is
\[
 \int_X(\langle\nabla h_a,\nabla h_b\rangle+2h_ah_b)dA
 +\sum_i\int_{C_i}
   (\langle\nabla h_a,\nabla h_b\rangle+
    \langle\nabla V_{i,a},\nabla V_{i,b}\rangle)dA,
\]
which is symmetric in $a,b$. This proves symmetry of $dF/ds$.
Since the positive orthant is simply connected, $F=dS$ for a function
$S$, and \eqref{A2b-positive-energy} makes $S$ strictly convex.
Finally, for distinct $s,t$, integrate along their straight segment:
\[
 (t-s)\cdot(F(t)-F(s))
 =\int_0^1 (t-s)^T(dF)_{s+v(t-s)}(t-s)\,dv>0.
\]
Thus $F$ is globally injective, as claimed.
\end{proof}

To specify the image, let $p\in[0,\infty)^I\setminus\{0\}$.
On $Y$ take the Jenkins--Strebel differential with cylinder circumferences
$2\pi\sqrt{p_i}$ at the points with $p_i>0$, all other marked points
being zero-weight marked points on its critical graph. Write $U_i(p)$
for its disc domains and $R_i(p)$ for their conformal radii. Define
\begin{equation}\label{A2b-support}
 K(p)=\sum_{p_i>0}p_i\log R_i(p),\qquad K(0)=0.
\end{equation}
Multiplying $p$ by a positive scalar rescales the differential and does
not change its disc domains. Equivalently, the weighted reduced-modulus
theorem gives
\begin{equation}\label{A2b-support-extremal}
 K(p)=\sup_{(D_i)}\sum_i p_i\log r_i(D_i),
\end{equation}
where the supremum ranges over pairwise disjoint pointed conformal
discs avoiding the other marked points; a zero-weight term is omitted.
Solynin's Theorem~2.1, equation (2.27), specializes to this assertion
\cite[\S1.2, (1.2), p.~9; \S2.3, (2.2), p.~23;
 Theorem 2.1, (2.27), p.~32]{Solynin1999A2b}:
take only the homotopy classes of circles about the positive-weight
points, no $G_j$ configurations in his equation (2.2), and use the
reduced disc modulus $(2\pi)^{-1}\log r_i$. His trajectory weights
$a_i=2\pi\sqrt{p_i}$ give the common factor $2\pi$ on both sides.
Zero-weight marked points are his allowed poles of order at most one;
such a pole may be removable when it meets a zero of the differential.
If one requires
nonzero competing discs at the zero-weight points too, the supremum
is unchanged: shrink the positive-weight extremal discs concentrically
by factors tending to $1$ and insert small discs at the remaining
points. This also
explains why the square roots of $p_i$, rather than $p_i$, are the
Strebel circumference weights.

The support value is computable directly from the unique differential,
without solving another hyperbolic capping problem. Write its principal
part at $x_i$ as $-p_i\,dz_i^2/z_i^2$. Choose a regular point $b_i$ of
the adjacent critical graph and the branch of
$\sqrt{-q_p/p_i}$ with residue $1$ at $x_i$. Then
\[
 \log R_i(p)=\lim_{z\to x_i}
 \left(\log|z_i(z)|-
  \operatorname{Re}\int_{b_i}^{z}\sqrt{-q_p/p_i}\right).
\]
The integral is taken inside the disc domain; its real part is
single-valued, and changing $b_i$ along the adjacent critical boundary
changes the integral only by an imaginary constant. This follows by
integrating $dw_i/w_i=\sqrt{-q_p/p_i}$ and normalizing $|w_i|=1$ on
that boundary.

\begin{lemma}[Proved: uniform cap bounds and recession]\label{A2b-recession}
There are constants $B_i<\infty$, depending only on the marked $Y$ and
the chosen tangent coordinates, such that every embedded pointed disc
avoiding the other marked points has conformal radius at most $B_i$.
For a canonical cap one has
\begin{equation}\label{A2b-collar-bound}
 r_i(s)\leq B_i\exp\left(
   -\frac{\arctan(\operatorname{csch}(\pi\sqrt{s_i}))}{\sqrt{s_i}}
                         \right).
\end{equation}
Moreover, for every $s_0>0$ and $p\geq0$,
\begin{equation}\label{A2b-recession-formula}
 \lim_{t\to\infty}\frac{S(s_0+tp)-S(s_0)}t=K(p).
\end{equation}
\end{lemma}
\begin{proof}
For the radius bound in genus zero, put the given point at $0$ and two
other marked points at $1,\infty$; Koebe's quarter theorem gives the
bound $4$, with the appropriate tangent-coordinate factor. In genus one,
lift the embedded disc to the universal covering flat plane. Its area
is at most the area $A$ of the torus, whereas the area formula for a
univalent map $f(z)=a_1z+\cdots$ gives area at least $\pi|a_1|^2$.
In higher genus, lift to the universal covering disc and apply
Schwarz's lemma. These arguments also cover extra marked points.

Double the hyperbolic surface across its geodesic boundaries. The collar
lemma gives an embedded one-sided collar of width
$w_i=\operatorname{arcsinh}(1/\sinh(L_i/2))$ on the original surface.
In the conformal coordinate above its width is
$a_i=\arctan(\sinh w_i)=\arctan(\operatorname{csch}(\pi\lambda_i))$.
The cap inverse map therefore extends univalently from the unit disc
to the disc of radius $e^{a_i/\lambda_i}$. The preceding global bound
applied to this enlarged disc proves \eqref{A2b-collar-bound}.
In particular $F_i(s)\to-\infty$ as $s_i\downarrow0$, uniformly in
the other parameters, and always $F_i(s)\leq\log B_i$.

The area comparison of Lemma~\ref{A2-strebel-deficit} gives the uniform
weighted deficit
\[
 0\leq K(s)-s\cdot F(s)\leq 2g-2+n.
\]
Indeed, every peripheral loop has grafted-metric length at least its
geodesic boundary length: project any flat-cylinder excursions to the
boundary and use the hyperbolic geodesic minimum. Cauchy--Schwarz on
each circle of the Strebel cylinders then compares their truncated
areas with the grafted area. The latter has hyperbolic part
$2\pi(2g-2+n)$ and the difference of its reduced flat areas is exactly
$2\pi s\cdot(F(s)-\log R(s))$.
Consequently, for $p>0$, the derivative
$\frac{d}{dt}S(tp)=p\cdot F(tp)$ tends to $K(p)$.
Ces\`aro averaging gives $S(tp)/t\to K(p)$.
This growth limit is unchanged by the fixed positive shift $s_0$:
the function $T(s)=S(s)-\sum_i(\log B_i)s_i$ is coordinatewise
nonincreasing, and, for $C$ such that $s_0\leq Cp$,
\[
 T((t+C)p)\leq T(s_0+tp)\leq T(tp).
\]
Dividing by $t$ proves \eqref{A2b-recession-formula} for $p>0$.
For a nonnegative $p$ with zero entries, feasibility of the cap discs
and \eqref{A2b-support-extremal} imply
$p\cdot F(s)\leq K(p)$, hence the upper bound $K(p)$ for the
recession limit. Set $b_i=\log B_i$. The coordinate derivative bounds
give, for $\epsilon>0$,
\[
 S(s_0+tp)\geq S(s_0+t(p+\epsilon\mathbf1))
                       -t\epsilon\sum_i b_i.
\]
The lower recession bound is therefore
$K(p+\epsilon\mathbf1)-\epsilon\sum_i b_i$.
From the supremum formula and the coordinate upper bounds,
$K(p+\epsilon\mathbf1)\to K(p)$: the upper bound is
$K(p)+\epsilon\sum b_i$, and any fixed nearly maximizing disc system
gives the required lower limit. Letting $\epsilon\downarrow0$ proves
the assertion, including the boundary directions.
\end{proof}

\begin{theorem}[Proved: exact image of the shadow map]\label{A2b-image}
The canonical cap radius map is a real-analytic diffeomorphism
\[
 (0,\infty)^I\longrightarrow\exp\mathcal L_Y,
 \qquad
 \mathcal L_Y=
 \left\{x\in\R^I:\ p\cdot x<K(p)
       \text{ for every }p\in[0,\infty)^I\setminus\{0\}\right\}.
\]
In particular its logarithmic image is open, convex, and downward closed.
It is precisely the set of conformal-radius vectors realized by systems
of pointed discs with pairwise disjoint analytic closures avoiding all
the other marked points. The boundary
support values are the weighted Strebel extrema \eqref{A2b-support}.
\end{theorem}
\begin{proof}
For a canonical system, one can enlarge any one cap a little into its
embedded hyperbolic collar while keeping every other cap fixed. For
any nonzero $p\geq0$, enlarge an index with $p_i>0$; the supremum formula
then gives the strict inequality $p\cdot F(s)<K(p)$.

Fix $x\in\mathcal L_Y$ and put $G(s)=S(s)-x\cdot s$.
Its sublevel sets are bounded. Indeed, otherwise choose $s_k$ with
$\|s_k\|\to\infty$ and $G(s_k)$ bounded above. From a fixed $s_0>0$,
pass to a subsequence such that
$p_k=(s_k-s_0)/\|s_k-s_0\|\to p\geq0$, $\|p\|=1$.
For any fixed $t>0$, convexity gives
\[
 G(s_0+tp_k)\leq
 (1-t/\|s_k-s_0\|)G(s_0)
 +(t/\|s_k-s_0\|)G(s_k).
\]
Taking the limit yields $G(s_0+tp)\leq G(s_0)$ for every $t$,
contradicting the positive recession slope $K(p)-p\cdot x$.

Extend $G$ to the closed orthant by lower semicontinuous closure.
A supporting affine function at any interior point bounds it below,
so this extension never takes the value $-\infty$. Bounded sublevel
sets imply that it attains its infimum. No minimum can lie at a finite
boundary point $b$. If the boundary value is infinite this is immediate.
Otherwise consider $g(t)=G(b+t\mathbf1)$ for $t>0$.
The convex closure has $g(t)\to G(b)$ as $t\downarrow0$; this follows
also by comparing with the segments from $b$ to an arbitrary interior
point. On the other hand,
\[
 g'(t)=\sum_i(F_i(b+t\mathbf1)-x_i)\longrightarrow-\infty:
\]
at least one $b_i$ is zero, its summand tends uniformly to $-\infty$
by \eqref{A2b-collar-bound}, and all other summands have finite uniform
upper bounds. Thus $g(t)<G(b)$ for small $t>0$, a contradiction.
The minimizer lies in the positive orthant and satisfies $F(s)=x$.
Strict convexity makes it unique. Positive definiteness of the Jacobian
then gives the real-analytic inverse.

The displayed half-space description proves convexity and downward
closure. Openness follows from the local inverse theorem (or from
continuity of $K$ on a normalized compact simplex). Every canonical
system has disjoint analytic closures. Conversely, a system with
disjoint analytic closures can be enlarged independently a little;
the same strict-inequality argument places its logarithmic radii in
$\mathcal L_Y$, so the result already proved supplies its unique
canonical representative with those radii. This last assertion
identifies radii, not the arbitrary discs with their canonical
representative.
\end{proof}

\paragraph{Dependency checkpoint.}
The moving-seam identity, its strictness, the collar bound, and the
convex minimization argument are proved here. The use of the
Strebel support values requires the weighted reduced-modulus theorem
in the indicated specialization. The recession calculation uses the
explicit area comparison of Lemma~\ref{A2-strebel-deficit}; it requires
no inference from convergence of abstract Teichm\"uller points to
convergence of embedded cap domains.
